\documentclass[11pt]{article}
\usepackage[margin=1in]{geometry}
\usepackage{amsmath,amssymb,amsthm,mathtools}
\usepackage{enumitem}
\usepackage{hyperref}
\hypersetup{colorlinks=true,linkcolor=blue,citecolor=blue,urlcolor=blue}

\newtheorem{theorem}{Theorem}
\newtheorem{proposition}{Proposition}
\newtheorem{lemma}{Lemma}
\newtheorem{corollary}{Corollary}
\newtheorem{definition}{Definition}
\newtheorem{hypothesis}{Hypothesis}
\newtheorem{remark}{Remark}
\newtheorem{warning}{Warning}
\newtheorem{conjecture}{Conjecture}
\newtheorem{problem}{Problem}

\newcommand{\Z}{\mathbb Z}
\newcommand{\Q}{\mathbb Q}
\newcommand{\F}{\mathbb F}
\newcommand{\C}{\mathbb C}
\newcommand{\G}{\Gamma}
\newcommand{\E}{\mathbb E}
\newcommand{\1}{\mathbf 1}
\newcommand{\ord}{\operatorname{ord}}
\newcommand{\lcm}{\operatorname{lcm}}
\newcommand{\im}{\operatorname{im}}
\newcommand{\rank}{\operatorname{rank}}
\newcommand{\Gal}{\operatorname{Gal}}

\title{A Kummer Distribution Criterion for Erd\H{o}s--Graham Problem \#203}
\author{Jared Wilder}
\date{Draft 0.1 --- 2026-05-12}

\begin{document}
\maketitle

\begin{abstract}
We formulate a precise Kummer distribution criterion sufficient for the negative resolution of Erd\H{o}s--Graham Problem \#203.  The criterion, called prime Hecke non-concentration (PHNC), asserts that fixed-radical Kummer characters do not concentrate almost entirely at the identity among primes \(q\le Q\), \(q\equiv1\bmod\ell\), uniformly for \(\ell\le(\log Q)^B\).  This note separates the exact algebraic reductions from the currently open analytic high-\(\ell\) seam.  It is a conditional criterion, not an unconditional solution of \#203.
\end{abstract}

\section{The problem}

For \((m,6)=1\), define
\[
A(m)=\{(k,\ell)\in\Z_{\ge0}^2:m2^k3^\ell+1\in\mathbb P\}.
\]
The desired negative resolution of Erd\H{o}s--Graham Problem \#203 is
\[
\forall m,\ (m,6)=1,\qquad |A(m)|=\infty.
\]

\section{Kummer vectors}

Let \(\ell\) be prime and \(q\equiv1\pmod\ell\).  Choose a primitive root \(g_q\bmod q\), and write
\[
2=g_q^{u_q},\qquad 3=g_q^{v_q}.
\]
Define
\[
x_\ell(q)=(u_q,v_q)\in\F_\ell^2.
\]

For \((r,s)\in\F_\ell^2\), define the Kummer character condition
\[
r u_q+s v_q\equiv0\pmod\ell,
\]
equivalently
\[
2^r3^s\in(\F_q^*)^\ell.
\]

\section{Prime Hecke non-concentration}

\begin{hypothesis}[PHNC]\label{H:phnc}
For every fixed \(B>0\), there exists \(C=C(B)>0\) such that, for every prime
\[
\ell\le(\log Q)^B
\]
and every nonzero \((r,s)\in\F_\ell^2\),
\[
\#\{q\le Q:q\equiv1\bmod\ell,\ 2^r3^s\in(\F_q^*)^\ell\}
\le
(1-\ell^{-C})\pi(Q;\ell,1)
\]
for all sufficiently large \(Q\).
\end{hypothesis}

\begin{remark}
PHNC is a non-concentration statement, not an equidistribution statement.  Equidistribution would predict relative density \(1/\ell\).  PHNC only forbids relative density \(1-o(\ell^{-C})\).
\end{remark}

\section{Kummer-BV criterion}

A stronger but more analytic form is a Kummer Bombieri--Vinogradov type variance theorem.

\begin{hypothesis}[Kummer-BV]\label{H:kbv}
For every \(A,B>0\), a suitable prime weight \(\omega_q\), and all \(Q\),
\[
\sum_{\ell\le(\log Q)^B}\ell^{\alpha-2}
\sum_{(r,s)\ne(0,0)}
\left|
\sum_{\substack{q\le Q\\ q\equiv1\bmod\ell}}
\omega_q e_\ell(r\log_q2+s\log_q3)
\right|^2
\ll_A(\log Q)^{-A}.
\]
\end{hypothesis}

\begin{remark}
The precise weight and exponent \(\alpha\) must be matched to the sieve application.  This draft records the structural role of the criterion.
\end{remark}

\section{PHNC to smoothing}

\begin{hypothesis}[Random-walk smoothing bridge]\label{H:krws}
PHNC implies a polynomial spectral gap for the Kummer-vector measures \(x_\ell(q)\), and after sufficiently many convolutions yields collision probability
\[
\ell^{-2}+O_A((\log Q)^{-A})
\]
uniformly in the required \(\ell\)-range.
\end{hypothesis}

\begin{remark}
This finite-field smoothing bridge is an explicit auxiliary lemma to prove.  One possible formulation is: if a measure on \(\F_\ell^2\) assigns at most \(1-\delta\) mass to every hyperplane, then its nontrivial Fourier coefficients are bounded away from \(1\) by a polynomial in \(\delta\), possibly under an anti-atom condition.
\end{remark}

\section{Conditional closure theorem}

\begin{theorem}[Conditional Erd\H{o}s--Graham \#203 criterion]\label{T:conditional}
Assume PHNC together with the random-walk smoothing bridge and the Kummer-BV-to-sieve implication.  Then for every \(m\in\Z_{>0}\) with \((m,6)=1\),
\[
|A(m)|=\infty.
\]
\end{theorem}

\begin{proof}[Proof architecture]
The implication chain is
\[
\mathrm{PHNC}
\Rightarrow
\mathrm{KRWS}
\Rightarrow
\mathrm{Kummer\ residue\ collision\ control}
\Rightarrow
\mathrm{rank\text{-}two\ Hecke\ large\ sieve}
\Rightarrow
\mathrm{Kummer\text{-}BV}
\Rightarrow
\mathrm{BVK\ sieve}
\Rightarrow
\mathrm{NEG\text{-}203}.
\]
The first arrow is Hypothesis \ref{H:krws}.  The Kummer residue collision control converts smoothed vector distribution into pair-collision bounds for the projective slopes.  The rank-two Hecke large sieve converts those collision bounds into the Kummer-BV variance estimate.  The BVK sieve then supplies infinitely many primes of the form \(m2^k3^\ell+1\).  Each arrow must be formalized as a separate proposition in a complete conditional paper.
\end{proof}

\section{Low-\(\ell\) range}

For a fixed nonzero \((r,s)\), let
\[
K_\ell=\Q(\zeta_\ell),
\qquad
L_{\ell,r,s}=K_\ell((2^r3^s)^{1/\ell}).
\]
By valuation at primes over \(2\) and \(3\), \(2^r3^s\) is not an \(\ell\)-th power in \(K_\ell\) when \((r,s)\ne0\).  Thus
\[
[L_{\ell,r,s}:K_\ell]=\ell.
\]
Effective Chebotarev/log-free Hecke estimates give PHNC in ranges where the field/conductor parameter is sufficiently small relative to \(Q\), heuristically
\[
\ell\log\ell\ll\log Q.
\]
A safe statement should cite the exact effective theorem and its field-invariant dependence.

\section{The high-\(\ell\) seam}

The open seam is
\[
(\log Q)^{1-\varepsilon}<\ell\le(\log Q)^B.
\]
Fixed-field ineffective Chebotarev is not enough to prove PHNC uniformly in this growing family.  The required theorem is the following.

\begin{problem}[Uniform fixed-radical Kummer non-concentration]
For every \(B>0\), prove that there exists \(C=C(B)>0\) such that PHNC holds for all
\[
\ell\le(\log Q)^B
\]
and all nonzero \((r,s)\in\F_\ell^2\).
\end{problem}

\section{GRH conditional variant}

\begin{hypothesis}[Kummer GRH]
Assume GRH for the Hecke/Artin \(L\)-functions attached to the fields
\[
\Q(\zeta_\ell,(2^r3^s)^{1/\ell})
\]
uniformly for \(\ell\le(\log Q)^B\) and nonzero \((r,s)\).
\end{hypothesis}

\begin{corollary}[GRH-conditional route]
Under a sufficiently uniform Kummer-GRH Chebotarev estimate, the PHNC criterion holds, and hence Theorem \ref{T:conditional} applies.
\end{corollary}

\begin{proof}
Uniform GRH Chebotarev gives strong splitting estimates in each cyclic Kummer layer.  These imply non-concentration at the identity value.  The conditional closure theorem then applies.
\end{proof}

\section{Conclusion}

This paper identifies a precise analytic target.  It does not claim that the target is currently proved in the high-\(\ell\) seam.  Its value is to make the route to \#203 exact.

\bibliographystyle{alpha}
\begin{thebibliography}{9}

\bibitem{LO}
J. C. Lagarias and A. M. Odlyzko.
\newblock Effective versions of the Chebotarev density theorem.
\newblock \emph{Algebraic Number Fields}, 1977.

\bibitem{TZ}
J. Thorner and A. Zaman.
\newblock A unified and improved Chebotarev density theorem.
\newblock \emph{Algebra \& Number Theory}, 13(5):1039--1068, 2019.

\bibitem{Washington}
L. Washington.
\newblock \emph{Introduction to Cyclotomic Fields}.
\newblock Springer.

\end{thebibliography}

\end{document}
